\documentclass[11pt]{article}
\usepackage[T1]{fontenc}
\usepackage[utf8]{inputenc}
\usepackage{lmodern,amsmath,amssymb,amsthm,mathtools}
\usepackage[margin=28mm]{geometry}
\usepackage{microtype,xcolor}
\usepackage[authoryear,round]{natbib}
\usepackage{xurl}
\usepackage[colorlinks=true,allcolors=blue!50!black]{hyperref}
\newtheorem{theorem}{Theorem}[section]
\newtheorem{proposition}[theorem]{Proposition}
\newtheorem{lemma}[theorem]{Lemma}
\newtheorem{corollary}[theorem]{Corollary}
\theoremstyle{definition}\newtheorem{definition}[theorem]{Definition}
\theoremstyle{remark}\newtheorem{remark}[theorem]{Remark}
\newcommand{\OO}{\mathcal O}
\newcommand{\PP}{\mathbf P}
\newcommand{\FF}{\mathbf F}
\newcommand{\Pic}{\operatorname{Pic}}
\newcommand{\NS}{\operatorname{NS}}
\newcommand{\End}{\operatorname{End}}
\newcommand{\Hom}{\operatorname{Hom}}
\newcommand{\Nm}{\operatorname{Nm}}
\newcommand{\Tr}{\operatorname{Tr}}
\newcommand{\ann}{\operatorname{ann}}
\newcommand{\rk}{\operatorname{rank}}
\newcommand{\diag}{\operatorname{diag}}
\newcommand{\workbenchurl}{https://github.com/KokunoYumeto/schemes-and-moduli-workbench}
\newcommand{\recordurl}{https://github.com/KokunoYumeto/mathematics-workbenches}
\title{Grothendieck's generic Frobenius ranks on products of elliptic curves, after Cartier and Tate}
\author{}
\date{23 November 2026}
\begin{document}
\maketitle
\begin{abstract}
Grothendieck's crystalline-cohomology notes ask how the generic stable Frobenius rank of hypersurface sections depends on their degree. We determine that rank in every positive degree for every cubed product polarization on a product of two elliptic curves in characteristic two. The answer depends only on the two elliptic $2$-ranks and on two parities. It includes every ordinary elliptic curve and every power of two in the polarization degrees. We also determine the rank in bidegree $(m,3)$ for every even $m\ge6$ when the second factor is supersingular. Cartier's operator gives the upper bounds, and Tate's residue duality supplies the matching lower bounds from explicit kernel--cokernel pairings. Cubic sections and nodal images under one-factor \'etale isogenies transfer these ranks to the required rectangular systems. A separate supersingular bundle calculation gives another uniform sparse family. These are complete linear-system results, with an explicit restriction-surjectivity qualification for ambient hypersurfaces. The mathematics and text were produced and checked by GPT-6 Astra Pro with the scopes stated at the end; no human has checked the mathematics.
\end{abstract}

\section{The question and the formula}

This paper computes the unit-root contribution of curves in fixed product linear systems. In his crystalline notes, Grothendieck asks about the maximal generic stable Frobenius rank of hypersurface sections and its dependence on the degree. The continuation of \S8.5.8 bears the marginal comment ``Or cela est \'etabli au n\textsuperscript{o}~7'' \citep[cote~4, archive sheets~7--8, PDF images~8--9]{Grothendieck}. The formula below gives an exact answer for cubed product polarizations in characteristic two. Its parity terms show why an ample degree alone does not determine the contribution of Frobenius: the two factor degrees interact differently with the ordinary and supersingular parts.

Let $k$ be an algebraically closed field of characteristic two. Let $E_1,E_2$ be elliptic curves with origins $O_1,O_2$, and put $A=E_1\times E_2$. Write $f_i\in\{0,1\}$ for the $2$-rank of $E_i$: thus $|E_i[2](k)|=2^{f_i}$, and $f_i=1$ means ordinary. The $2$-rank $f_C$ of a smooth projective curve is that of its Jacobian. For positive integers $m,n$ write
\[
 N(m,n)=\OO_{E_1}(mO_1)\boxtimes\OO_{E_2}(nO_2).
\]
If $C\in|N(m,n)|$ is smooth, define its relative Frobenius space and rank by
\[
 Q_C=\operatorname{coker}\bigl[H^1(A,\OO_A)\longrightarrow H^1(C,\OO_C)\bigr],
 \qquad r(C)=\dim_k\bigcap_{j\ge0}\operatorname{im}(F^j|_{Q_C}).
\]
Here $F$ is absolute Frobenius, so it is semilinear. The indicated restriction is injective. We prove below that
\begin{equation}\label{eq:dim}
 g(C)=mn+1,\qquad \dim Q_C=mn-1,\qquad r(C)=f_C-f_1-f_2.
\end{equation}
The generic value means the value on a nonempty open of the smooth-member locus of the complete linear system. Put $\epsilon(t)=1$ for even $t$ and $\epsilon(t)=0$ for odd $t$.

\begin{theorem}[All rectangular product polarizations]\label{thm:main}
For any positive integers $a,b,d$, a general smooth curve in $|N(3ad,3bd)|$ has relative stable Frobenius rank
\begin{equation}\label{eq:main}
 r_d=9abd^2-1-(1-f_2)\epsilon(ad)-(1-f_1)\epsilon(bd).
\end{equation}
Its genus is $9abd^2+1$ and its relative coherent dimension is $9abd^2-1$. There is no restriction on the ordinary elliptic parameter or on the $2$-adic valuations of $a,b,d$.

Over a perfect field the same assertion holds for geometric generic rank after algebraic closure. For the embedding defined by $L=N(3a,3b)$, formula \eqref{eq:main} holds for ambient degree-$d$ hypersurfaces whenever restriction onto $H^0(A,L^d)$ is surjective; a sufficient range is
\[
 d\ge\max\{1,\operatorname{reg}(\mathcal I_A)-1\}.
\]
\end{theorem}

In words, two ordinary factors give full relative rank. With an ordinary first factor and a supersingular second factor, exactly one rank is lost when the first degree is even. With two supersingular factors, each even degree loses one rank. These are ranks in a fixed complete system; they are not generic assertions in the moduli space of all curves. For example, in bidegree $(6,3)$ the possibilities are
\[
\begin{array}{c|c|c|c}
(f_1,f_2)&g&f_C&r(C)\\\hline
(1,1)&19&19&17\\
(1,0)&19&17&16\\
(0,1)&19&18&17\\
(0,0)&19&16&16.
\end{array}
\]
Changing which factor is ordinary changes the answer even though the ambient $2$-rank remains one.

The difficult small systems have one factor of degree three and an even degree on the other. The two uniform inputs are useful independently.

\begin{theorem}[The even thin systems]\label{thm:thin}
Let $E_{\rm o}$ be any ordinary elliptic curve and $E_{\rm s}$ any supersingular elliptic curve over $k$.
\begin{enumerate}
\item For every even $m\ge6$, a general member of
$|\OO_{E_{\rm o}}(mO)\boxtimes\OO_{E_{\rm s}}(3O)|$ has genus $3m+1$, full $2$-rank $3m-1$, and relative rank $3m-2$.
\item For every even $m\ge6$, a general member of
$|\OO_{E_{\rm s}}(mO)\boxtimes\OO_{E_{\rm s}}(3O)|$ has genus $3m+1$ and full and relative rank $3m-2$. Proposition~\ref{prop:ss-all-even} gives an explicit section in every such degree. For $m=2q$ with $6\mid q$, Proposition~\ref{prop:ss-sparse} gives a second explicit family with first and second Cartier ranks $6q-2$.
\end{enumerate}
\end{theorem}

\paragraph{Sources and context.}
The basic Cartier isomorphism and its product and local formulas are classical: see \citet[0, (2.1.4)--(2.1.13), pp.~514--515]{Illusie}. The link with the $p$-rank of a Jacobian is part of the coherent and Dieudonn\'e description in \citet[Propositions~3.18 and 3.22--3.24, pp.~89, 99--101]{Oda}. The ordinary Frobenius splitting is also a rank-two case of \citet[Corollary~1.7, pp.~53--54]{Oda1971}; we prove the splitting and the additional endomorphism calculations directly. The residue proof uses \citet[Theorems~2--3 and the corollary, pp.~155--156; \S4 and Theorem~5, pp.~157--159]{Tate}; the finite-isogeny inputs are stated at their precise locations in Section~\ref{sec:norm}. These classical results carry the proof. We prove the fixed-system rank formula, the uniform ordinary-parameter determinant, and the supersingular two-point bundle calculation in the forms needed here.

For neighboring moduli questions, \citet[Theorem~1.1, p.~2]{ChangEtAl} determine the dimensions of $p$-rank strata of double covers of a fixed elliptic curve in characteristic $p>2$. Their variable covering space and characteristic hypothesis differ from the fixed product systems considered here. Recent AI-assisted work of \citet[Introduction, Theorems~A--C, and acknowledgements]{Zou} relates quasi-$F$-splitting and crystalline slopes for higher-dimensional symplectic varieties. It provides context for Frobenius questions, but no input to this proof. A bounded search of journal, author, arXiv, ECCC, and Hexagon sources did not identify the exact rectangular formula; no priority claim is made.

\paragraph{Why the proof works.}
A supersingular factor kills the top invariant form under Cartier. An even factor degree allows a global Frobenius connection, and differentiating the defining section then supplies another Cartier-kernel vector. These give the parity upper bounds. A single cubic section starts the odd cases. For the ordinary thin systems, the Frobenius algebra splits into two degree-zero lines. Two pencils turn stable rank into a $2\times2$ residue determinant, and evaluation on reduced divisors proves that determinant nonzero for every ordinary elliptic curve. For every even supersingular degree, an explicit section has kernel bundle $F_*\OO_E\oplus\OO_E$; two residue functionals give an invertible kernel--cokernel pairing. Another sparse family is treated by the nonsplit Frobenius bundle and its two rank-drop points. Finally, \'etale isogenies acting in one factor produce nodal images without changing the curve defect, and a transverse union handles the two-even supersingular case. This also explains where characteristic two, product structure, and the factor degrees enter.

\section{Cohomology, Cartier, and the parity bounds}\label{sec:cohom}

We first relate the desired relative rank to a finite Cartier matrix, and explain why a singular example can prove a generic smooth lower bound.

\begin{lemma}[Stable parts and divisor cohomology]\label{lem:cohom}
For $m,n>0$, every nonzero divisor $C\in|N(m,n)|$ is connected and has $h^1(C,\OO_C)=mn+1$. There is an exact Frobenius-compatible sequence
\begin{equation}\label{eq:divseq}
0\to H^1(A,\OO_A)\to H^1(C,\OO_C)\to
H^2(A,N^{-1})\to H^2(A,\OO_A)\to0.
\end{equation}
The endomorphism of $H^2(A,N^{-1})$ in this sequence is $sF$, where $s$ defines $C$: Frobenius first maps to $H^2(A,N^{-2})$ and multiplication by $s$ maps back. Its dual is multiplication by $s$ followed by Cartier trace. The relative stable rank is
\[
 r(C)=\rk_{\rm st}(sF|H^2(A,N^{-1}))-f_1f_2.
\]
If $C$ is smooth, \eqref{eq:dim} holds.
\end{lemma}
\begin{proof}
K\"unneth and the elementary cohomology of a positive-degree line bundle on an elliptic curve give $H^i(A,N)=0$ for $i>0$, $h^0(N)=mn$, $H^1(A,N^{-1})=0$, and $h^2(N^{-1})=mn$. Apply the divisor sequence $0\to N^{-1}\xrightarrow{s}\OO_A\to\OO_C\to0$. It gives $H^0(\OO_C)=k$, connectedness, \eqref{eq:divseq}, and $h^1(\OO_C)=2+mn-1=mn+1$. For a smooth curve this is also adjunction, since $K_A=0$ and $N^2=2mn$.

For a semilinear map on a finite-dimensional vector space over a perfect field, the descending images stabilize. On their intersection the map is bijective, and on the complementary stabilized kernel it is nilpotent. This follows by taking a power for which both image and kernel dimensions have stabilized, and applying rank--nullity to that power. The stable-part functor is exact. Indeed a lift of a vector in a stable quotient can first be projected to the stable summand of the middle term; its nilpotent summand has nilpotent image and therefore zero stable projection. Kernels are treated by the same decomposition.

The compatible map on $N^{-1}$ is $sF$: for a local section $v$, the two routes to $\OO_A$ both give $s^2v^2$. Duality identifies its transpose with $u\mapsto C(su\omega)/\omega$, where $\omega$ is a product of invariant differentials. This is also checked directly from the local Cartier coefficient formula. On $H^1(A,\OO_A)$ the stable dimension is $f_1+f_2$, and on $H^2(A,\OO_A)$ it is $f_1f_2$, by K\"unneth and the elliptic Cartier criterion. Taking stable parts proves the assertion. For smooth $C$, the classical Cartier--Jacobian identification says that the stable dimension of $H^1(C,\OO_C)$ is $f_C$, giving \eqref{eq:dim}.
\end{proof}

Fixing a basis identifies $sF$ with $v\mapsto M(s)v^{[2]}$. Its $j$th iterate has matrix
\[
 M(s)M(s)^{[2]}\cdots M(s)^{[2^{j-1}]},
\]
where brackets denote entrywise powers. In dimension $mn$, the stable rank equals the rank of the $mn$th iterate. Its rank-at-least loci are open because its entries are polynomial functions of the coefficients of $s$. Thus the generic stable rank is maximal. When $m,n\ge3$, the exterior product of the two very ample elliptic line bundles is very ample. The nonempty smooth-member open therefore meets every nonempty rank open.

We use absolute Cartier as a $2^{-1}$-linear operator. If its columns in a fixed basis are the columns of $T$, so that it sends $v$ to $T v^{[1/2]}$, coherent Frobenius has matrix
\begin{equation}\label{eq:twist}
 M=(T^{[2]})^{t},\qquad F^2\text{ has matrix }MM^{[2]}.
\end{equation}
All comparisons between two iterates below use this convention. In a model and section defined over $\FF_2$, the entries of $T$ lie in $\FF_2$ and ordinary matrix powers compute the iterate ranks.

\begin{lemma}[Parity bounds]\label{lem:upper}
Suppose $m,n\ge3$ and at least one factor is supersingular. In the complete system $|N(m,n)|$, the generic relative stable rank is at most
\[
 mn-1-(1-f_2)\epsilon(m)-(1-f_1)\epsilon(n).
\]
\end{lemma}
\begin{proof}
Let $\omega_i$ be invariant differentials and let $\partial_i$ be the dual vector fields. Since $C(\omega_1\wedge\omega_2)=0$, multiplication--Cartier kills $s$ itself:
$C(s^2\omega_1\wedge\omega_2)=sC(\omega_1\wedge\omega_2)=0$.
If $m$ is even, the first factor of $N$ is a Frobenius pullback. Its horizontal square transition functions define a global action of $\partial_1$ on $H^0(N)$. If $E_2$ is supersingular, then in such a trivialization, writing $s=f$,
\[
 f\partial_1f\,\omega_1\wedge\omega_2=f\,df\wedge\omega_2.
\]
The one-form $f\,df$ is closed in characteristic two. Cartier is multiplicative on closed forms and $C(\omega_2)=0$, so $C_s(\partial_1s)=0$. Transitions depending only on the second coordinate do not affect this global identity. The same argument gives $\partial_2s$ in the kernel when $n$ is even and $E_1$ is supersingular.

These kernel sections are generically independent. In the supersingular model $y^2+y=x^3$, the invariant derivation has $\partial x=1$, $\partial y=x^2$ and $\partial^2=0$. A relation $\partial s=cs$ forces $c=0$, and the derivation is nonzero on the section $x$. For an ordinary curve one can use $y^2+xy=x^3+c$, $c\ne0$, and $\partial x=x$, $\partial y=x^2+y$, $\partial^2=\partial$. The sections $1+x$ and $x$ are independent and witness independence of $s,\partial s$. These rational sections lie in the indicated spaces because an even degree at least three is at least four. If both factors are supersingular and both degrees are even, $s=x_1x_2$ gives the independent triple $x_1x_2,x_2,x_1$. Independence is open in the section space. The associated larger minors vanish on a dense open, hence vanish identically. First rank bounds stable rank, and Lemma~\ref{lem:cohom}, with $f_1f_2=0$, gives the result.
\end{proof}

\begin{lemma}[Nodal defect]\label{lem:nodal}
Let $C$ be a connected reduced projective nodal curve, with smooth normalization components $C_i$, and let $b$ be the first Betti number of its dual graph. Then
\[
 h^1(C,\OO_C)=\sum_i g(C_i)+b,\qquad
 \rk_{\rm st}H^1(C,\OO_C)=\sum_i f_{C_i}+b.
\]
Consequently its defect, the difference of these two integers, is $\sum_i(g(C_i)-f_{C_i})$.
\end{lemma}
\begin{proof}
In the normalization sequence, the quotient of $\OO_C$ in $\bigoplus_i\OO_{C_i}$ is one copy of $k$ for each node. The map on constants is the graph incidence map. Its cokernel is $k^b$, and the resulting exact sequence is
$0\to k^b\to H^1(C,\OO_C)\to\bigoplus_iH^1(C_i,\OO_{C_i})\to0$.
Frobenius acts by coefficientwise squaring on the graph term and is bijective. Exactness of stable parts proves both formulas.
\end{proof}

\section{A cubic calculation on the two elliptic models}\label{sec:cubic}

We now give a small section that starts all odd-degree constructions, keeping the parameter of the ordinary curve arbitrary.

Over $k$, every ordinary elliptic curve has a model
\[
 E_c:\ y^2+xy=x^3+c,\quad c\ne0,\qquad \omega=dx/x,
\]
and every supersingular elliptic curve is isomorphic to
\[
 E_{\rm s}:\ y^2+y=x^3,\qquad \omega=dx.
\]
The standard Weierstrass criterion is the nonvanishing of $a_1$ in characteristic two; see \citet[Example~5.26, p.~82]{EGM}. Here is the required normalization. In the ordinary case scale to $a_1=1$, translate $x$ to remove $a_3$, and translate $y$ by a constant to remove $a_4$. A further replacement $y\mapsto y+ux$ removes $a_2$ by solving $u^2+u=a_2$. Nonsingularity then says $c\ne0$. In the supersingular case scale $a_3$ to one. In $y^2+y=x^3+a_2x^2+a_4x+a_6$, replacements $x=X+r$, $y=Y+sX+t$ remove $X^2$ and $X$ by taking $s=r^2+a_4$ and a root of $r^4+r+a_4^2+a_2=0$; a solution of an Artin--Schreier equation removes the constant. Algebraic closedness provides these roots.

In the basis $(1,x,y)$ of $H^0(\OO(3O))$, the two multiplication--Cartier pairings are
\begin{equation}\label{eq:beta3}
 \beta_{\rm o}=
 \begin{pmatrix}1&0&x\\0&x&y+\sqrt c\\x&y+\sqrt c&y\end{pmatrix},
 \qquad
 \beta_{\rm s}=
 \begin{pmatrix}0&1&x\\1&0&y\\x&y&0\end{pmatrix}.
\end{equation}
For example, writing a rational function as $a^2+xb^2$, Cartier with respect to $dx$ returns $b$, and Cartier with respect to $dx/x$ returns $a$. On $E_c$, put $u=y+\sqrt c$. The identity $u^2=xu+x\sqrt c+x^3$ is equivalently $y^2=xy+x^3+c$, and $y=x^2+x((y+\sqrt c)/x)^2$. Applying these coefficient formulas to the six products proves \eqref{eq:beta3}; on $E_{\rm s}$ use $y=y^2+x^3$.

\begin{proposition}[The same cubic section in both cases]\label{prop:cubic}
For $E_1=E_c$ or $E_{\rm s}$ and $E_2=E_{\rm s}$, the section
\[
                         s=y+xz+w
\]
of $N(3,3)$ has Cartier stable rank eight. Consequently a general smooth member has defect one in the mixed case and defect two in the supersingular case.
\end{proposition}
\begin{proof}
Order the tensor basis by $3i+j$, with both factor bases $(1,x,y)$ and $(1,z,w)$ and indices starting at zero. Tensor the pairings \eqref{eq:beta3}, contract with the three nonzero coefficients of $s$, and use \eqref{eq:twist}. For two supersingular factors, the minor of $MM^{[2]}=M^2$ obtained by removing row and column two is a permutation matrix: its row nonzero positions, in the remaining order, are
\[
                        1,2,3,6,5,4,7,8.
\]
Its determinant is one. For the mixed case remove row two and column zero. The resulting minor is
\begin{equation}\label{eq:cubic-minor}
 \begin{pmatrix}
0&1&1&0&0&0&0&0\\
0&0&1&1&0&0&0&0\\
0&0&1&0&0&0&0&0\\
c&0&0&1&0&1&0&0\\
c^2&c&c&0&0&1&1&0\\
c&0&0&1&0&0&0&0\\
c^2&0&0&1&1&1&1&0\\
0&c^2&0&c&1&1&1&1
 \end{pmatrix}.
\end{equation}
Expand in the last column and use the first three rows to eliminate columns two, three and four. Expansion in the fifth column then leaves
$\left(\begin{smallmatrix}c&1&0\\c^2&1&1\\c&0&0\end{smallmatrix}\right)$, of determinant $c$. Thus the second coherent Frobenius rank is at least eight in both cases. The first rank is at most eight by Lemma~\ref{lem:upper}, so equality of the first two ranks gives stable rank eight. The rank open meets the smooth open. In the mixed case the ambient stable rank is one, giving $f_C=9$ and genus ten; in the supersingular case it is zero, giving $f_C=8$.
\end{proof}

\section{All even thin systems with an ordinary factor}\label{sec:mixed}

This section proves the first part of Theorem~\ref{thm:thin} by a residue determinant, uniformly in the ordinary elliptic curve and the degree.

\begin{lemma}[The ordinary Frobenius algebra]\label{lem:ordinary-algebra}
Let $F:E\to E'=E^{(2)}$ be relative Frobenius on an ordinary elliptic curve, and let $H=\OO_{E'}(qO)$ with $q\ge3$. There is an orthogonal decomposition
\[
 F_*\OO_E=\OO_{E'}\oplus B,\qquad B\not\simeq\OO_{E'},\qquad
 \mu:B^2\xrightarrow{\sim}\OO_{E'},
\]
for which multiplication followed by normalized Cartier is
$\langle(a,b),(c,d)\rangle=ac+\mu bd$.
Put $V=F_*F^*H=H\oplus HB$. One can choose sections $a_0\in H^0(H)$ and $b_2\in H^0(HB)$ with reduced disjoint zero divisors $D_0,D_2$. For any such choice all three maps
\begin{align*}
 H^0(H)\oplus H^0(HB)&\longrightarrow H^0(H^2),
 &(x,y)&\longmapsto a_0x+\mu b_2y,\\
 H^0(HB)&\longrightarrow H^0(D_0,HB|_{D_0}),&&\\
 H^0(H)&\longrightarrow H^0(D_2,H|_{D_2})&&
\end{align*}
are isomorphisms.
\end{lemma}
\begin{proof}
Choose invariant differentials on $E,E'$ so that Cartier gives an $\OO_{E'}$-linear functional $\tau:F_*\OO_E\to\OO_{E'}$ with $\tau(1)=1$. This is possible precisely because $E$ is ordinary. The pairing $\tau(uv)$ is perfect by the local two-basis formula. Its orthogonal complement to the unit is a line bundle $B=\ker\tau$. The square of a local generator of $B$ lies in the unit summand, since relative Frobenius identifies the functions on the base with squares. Perfection of the pairing makes this square a unit, giving $\mu$. Since $H^0(F_*\OO_E)=k$, one has $H^0(B)=0$; a degree-zero line bundle with this property is nontrivial and has $H^1(B)=0$ as well.

Both degree-$q$ bundles $H,HB$ are very ample. Their general sections have reduced zero divisors, and a general second section avoids the finite zero divisor of the first. The exact sequence
\[
0\longrightarrow B\xrightarrow{t\mapsto(\mu b_2t,a_0t)}
 H\oplus HB\xrightarrow{(x,y)\mapsto a_0x+\mu b_2y}H^2\longrightarrow0
\]
is checked locally wherever $a_0$ or $b_2$ is a unit. The acyclicity of $B$ gives the first isomorphism. The evaluation sequences have kernels $HB(-D_0)=B$ and $H(-D_2)=B^{-1}=B$, proving the other two.
\end{proof}

\begin{proposition}[The residue determinant]\label{prop:residue}
There are $a_1\in H^0(H)$ and $b_1\in H^0(HB)$ such that, for
\[
 f_0=(a_0,0),\qquad f_1=(a_1,b_1),\qquad f_2=(0,b_2)
 \quad\text{in }H^0(V),
\]
the section $s=f_0+zf_1+wf_2$ on $E\times E_{\rm s}$ has Cartier first and stable rank $6q-2$ in bidegree $(2q,3)$.
\end{proposition}
\begin{proof}
First require $(a_0,a_1)$ and $(b_1,b_2)$ to be basepoint-free pencils. This is the nonempty open defined by $a_1$ nonzero on $D_0$ and $b_1$ nonzero on $D_2$. Write an input triple as $(x_i,y_i)$ with $x_i\in H^0(H)$ and $y_i\in H^0(HB)$. Tensoring the orthogonal ordinary pairing with $\beta_{\rm s}$ gives the linearized Cartier equations
\begin{align}
 a_0x_1+a_1x_0+\mu b_1y_0&=g_0,\label{eq:res0}\\
 a_0x_2+\mu b_2y_0&=g_1,\label{eq:res1}\\
 a_1x_2+\mu b_1y_2+\mu b_2y_1&=g_2.\label{eq:res2}
\end{align}
The middle equation determines $x_2,y_0$ uniquely. The remaining two maps are pencil multiplication maps, each with one-dimensional kernel and cokernel, as follows from
$0\to\OO\to H^{\oplus2}\to H^2\to0$ and the corresponding sequence for $HB$. Consequently $T_s$ has rank $6q-2$ and kernel spanned by
\[
 k_a=((a_0,0),(a_1,0),(0,0)),\qquad
 k_b=((0,0),(0,b_1),(0,b_2)).
\]

To test stable rank one must compare the correctly twisted kernel and cokernel. Squaring and descent along Frobenius give a bijective Frobenius-semilinear map
\[
 Q:H^0(E,F^*H)\longrightarrow H^0(E',H^2),
 \qquad Q(a,b)=a^2+\mu b^2.
\]
This notation expresses descent of the square of a section; it is not ordinary multiplication of sections of $H$ and $HB$ before descent. For origin-divisor bundles the target is canonically $H^2$. The absolute Cartier map is $Q^{-1}T_s$, componentwise on the second-factor basis. Thus the kernel meets the image trivially exactly when
$Q(\ker T_s)\to\operatorname{coker}T_s$ is invertible. The two vectors to test are
\[
 Q(k_a)=(a_0^2,a_1^2,0),\qquad
 Q(k_b)=(0,\mu b_1^2,\mu b_2^2).
\]

Fix an invariant differential $\omega$ on $E'$. Residue duality identifies the two pencil cokernels with $k$ through
\[
 L_0(g)=\sum_{P\in D_0}\operatorname{Res}_P\frac{g\omega}{a_0a_1},
 \qquad
 L_2(g)=\sum_{P\in D_2}\operatorname{Res}_P\frac{g\omega}{\mu b_1b_2}.
\]
Indeed on the opens where the two respective pencil sections are nonzero, the two local solutions differ by the indicated fraction times the kernel generator. Their difference is the \v Cech connecting cocycle. Its pairing with $\omega$ is the sum of these residues. The local coefficient rule and this duality are \citet[Theorem~2, p.~155; \S4, pp.~157--159]{Tate}.

Reducing \eqref{eq:res1} at $D_0$ gives $y_0=g_1/(\mu b_2)$, and reducing it at $D_2$ gives $x_2=g_1/a_0$. The discarded differences in \eqref{eq:res0} and \eqref{eq:res2}, after division by the respective pencil products, are regular at the residue points: they are $b_1x_2/(b_2a_1)$ and $a_1y_0/(a_0b_1)$. Applying the two functionals to $Q(k_a),Q(k_b)$ therefore gives
\begin{equation}\label{eq:resmatrix}
 \begin{pmatrix}t&u\\v&t\end{pmatrix},
\end{equation}
where
\begin{align*}
 t&=\sum_{D_0}\operatorname{Res}\frac{a_1b_1}{a_0b_2}\omega,\\
 u&=\sum_{D_0}\operatorname{Res}\frac{\mu b_1^3}{a_0a_1b_2}\omega,\\
 v&=\sum_{D_2}\operatorname{Res}\frac{a_1^3}{\mu a_0b_1b_2}\omega.
\end{align*}
The lower-right entry is the expression for $t$ summed over $D_2$. That differential has poles only in $D_0\cup D_2$, so the global residue theorem and characteristic two make the two sums equal. Every displayed quotient is a rational function: $\mu$ cancels the $B$-factors.

Replace $b_1$ by $b_1+\lambda b_2$. The pencil stays basepoint-free because its values on $D_2$ stay fixed. The entry $v$ is constant since only those values enter its simple-pole residues. The change of $t$ is
$\lambda\sum_{D_0}\operatorname{Res}(a_1\omega/a_0)=0$ by the global residue theorem, since this differential has no other poles. The entry $u$ is a polynomial of degree at most three in $\lambda$, with linear coefficient
\begin{equation}\label{eq:c1}
 c_1=\sum_{D_0}\operatorname{Res}\frac{\mu b_1^2}{a_0a_1}\omega.
\end{equation}
For each fixed allowed $a_1$, this is a nonzero quadratic polynomial in $b_1$. The evaluation isomorphism of Lemma~\ref{lem:ordinary-algebra} gives independent coordinates $b_1(P)$ at the reduced points of $D_0$. In those coordinates \eqref{eq:c1} is a sum of their squares with nonzero coefficients, because $a_0$ has simple zeros and $a_1,\mu,\omega$ are units there.

Similarly, for each fixed allowed $b_1$, $v$ is a nonzero cubic polynomial in $a_1$: evaluation at $D_2$ is an isomorphism and every simple-pole coefficient is nonzero. The allowed pairs form an irreducible nonempty open in a vector space. The two additional nonvanishing opens $c_1\ne0$ and $v\ne0$ therefore meet. Choose a pair in their intersection. Then $t$ and $v$ are constant under the shift, and $u(\lambda)$ is nonconstant. The determinant $t^2+u(\lambda)v$ is a nonzero polynomial, so it is nonzero for some $\lambda\in k$. This proves that \eqref{eq:resmatrix} is invertible for a choice of sections. Kernel and image meet trivially, so first rank already equals stable rank, namely $6q-2$.
\end{proof}

\begin{proof}[Proof of Theorem~\ref{thm:thin}(1)]
Write $m=2q$. Proposition~\ref{prop:residue} gives a section with stable rank $3m-2$. Its defining curve need not be smooth: the nonzero iterated-rank minor defines a nonempty open in the complete system. This open meets the smooth open because both degrees are at least three. Lemma~\ref{lem:upper} is the matching upper bound. Lemma~\ref{lem:cohom} then adds the ambient rank one and gives genus $3m+1$ and full rank $3m-1$.
\end{proof}

\section{The supersingular Frobenius bundle}\label{sec:ssbundle}

The ordinary splitting is unavailable on a supersingular elliptic curve. We replace it by an explicit description of the endomorphisms of its nonsplit Frobenius bundle.

Throughout this section let $E/\FF_2$ be $y^2+y=x^3$, let $O$ be its point at infinity, and identify its Frobenius twist with $E$ using this equation. Write $F$ for relative Frobenius and $\partial$ for the invariant derivation $\partial x=1$, $\partial y=x^2$.

\begin{lemma}[A rank-two bundle and its endomorphisms]\label{lem:ss-end}
For $q>0$, put $H=\OO_E(qO)$ and $V=F_*\OO_E(2qO)$. Then $V=U\otimes H$, where $U=F_*\OO_E$ is a nonsplit extension of $\OO_E$ by itself. Multiplication followed by Cartier with respect to $dx$ gives a perfect alternating pairing $V\otimes V\to\det V=H^2$. Furthermore,
\[
 h^0(\End V)=2,\qquad h^0(\End V/\OO_E\operatorname{id})=2.
\]
The global endomorphisms are spanned by the identity and $\partial$, both of trace zero. The global classes modulo scalars are spanned by multiplication by the source function $x$ and by $\partial$.
\end{lemma}
\begin{proof}
Cartier kills $dx$, so Frobenius is zero on $H^1(E,\OO_E)$. In
$0\to\OO_E\to F_*\OO_E\to B^1_E\to0$, the quotient is a line bundle of degree zero. The cohomology sequence gives it a nonzero global section, hence it is $\OO_E$. The extension is nonsplit because $h^0(F_*\OO_E)=1$. Projection formula gives $V=U\otimes H$ and $\det V=H^2$. In the rational basis $(1,x)$, Cartier satisfies $C(dx)=C(x^2dx)=0$, $C(xdx)=dx$, so its pairing is the alternating wedge pairing. The local Cartier transformation law makes it global and perfect.

Dualizing the nonsplit extension again gives a nonsplit extension. Thus $\dim\Hom(U,\OO_E)=1$. In the sequence obtained by applying $\Hom(U,-)$ to the original extension, the identity maps to its nonzero quotient map $U\to\OO_E$. Therefore $\dim\End(U)=2$. The identity and the nonzero nilpotent endomorphism $\partial$ are independent and span it. They both have trace zero in characteristic two.

The trace pairing identifies $\End(U)$ with its dual. The map $H^1(\OO_E)\to H^1(\End U)$ from the scalar inclusion is therefore dual to the zero global trace map. The cohomology sequence of the quotient by scalars gives dimension two for its global sections.

The derivative acts globally on $V$ because $\OO_E(2qO)=F^*H$ has horizontal square transitions. Multiplication by $x$ defines a global class modulo scalars, although $x$ has a pole at $O$. Indeed, for the uniformizer $t=x/y$,
\begin{equation}\label{eq:at-infty}
                       x+t^{-2}=y/x^2.
\end{equation}
The right side is regular at $O$, while $t^{-2}$ is a square and therefore a scalar over the Frobenius base. Away from $O$, multiplication by $x$ is regular. In the rational basis $(1,x)$, denoting the base coordinate by $X$, the two classes have representatives
\begin{equation}\label{eq:ss-endmat}
 \begin{pmatrix}0&X\\1&0\end{pmatrix},\qquad
 \begin{pmatrix}0&1\\0&0\end{pmatrix}.
\end{equation}
They are independent modulo scalars, so the dimension already computed shows that they span. This proves the required bundle facts directly.
\end{proof}

\begin{proposition}[A uniform sparse section]\label{prop:ss-sparse}
For every positive $q$ divisible by six, define a section of $N(2q,3)$ on $E\times E$ by
\begin{align*}
 s_q&=A+Bz+Cw,\qquad P=x^{q-3},\\
 A&=x+P\{x^2+(1+x)y\},\qquad B=x+Py^2,\qquad C=1+Py.
\end{align*}
Its first and second Cartier ranks are both $6q-2$. A general smooth member of $|N(2q,3)|$ therefore has full and relative stable rank $6q-2$.
\end{proposition}
\begin{proof}
The pole orders of $A,B,C$ are $2q-1,2q,2q-3$, respectively. They are thus sections of $\OO_E(2qO)$. Put $u=x^{q/2-2}$, so $P=xu^2$. Reducing by $y^2+y=x^3$ gives the two-basis decomposition
\begin{equation}\label{eq:ss-ab}
 (A,B,C)=a^2+xb^2,
 \quad
 \begin{cases}
 a=(xu(x+y),\ 0,\ 1+x^2u),\\
 b=(1+u(x+y+x^2),\ 1+uy,\ uy).
 \end{cases}
\end{equation}
For instance, the last component is
$(1+x^2u)^2+x(uy)^2=1+x^4u^2+xu^2(y+x^3)=1+xu^2y=C$;
the first two components follow by the same substitution. Because all coefficients are in $\FF_2$, $a,b$ are precisely the coordinate rows of the three sections $f_0=A,f_1=B,f_2=C$ of $V=F_*\OO_E(2qO)$ on the identified Frobenius base.

Consider the map
\begin{align*}
 \Psi:V^3&\longrightarrow(\det V)^3,\\
 (v_0,v_1,v_2)&\longmapsto
 (f_0\wedge v_1+f_1\wedge v_0,\
  f_0\wedge v_2+f_2\wedge v_0,\
  f_1\wedge v_2+f_2\wedge v_1).
\end{align*}
The second-factor pairing \eqref{eq:beta3} says exactly that $H^0(\Psi)$ is multiplication--Cartier, in the ordered tensor basis. We determine its kernel and cokernel as sheaves.

First determine where the three $f_i$ fail to span $V$. Put $\beta=1+uy$. On $\beta\ne0$, the two rows in \eqref{eq:ss-ab} have rank one if and only if
\[
                     xu(x+y)=0,\qquad 1+x^2u=0.
\]
The second equation forces $x\ne0$; the first then gives $y=x$. The elliptic equation becomes $x^2+x+1=0$. Conversely $q/2$ is divisible by three, so both points
\[
 Q_\zeta=(\zeta,\zeta),\qquad Q_{\zeta^2}=(\zeta^2,\zeta^2),
 \qquad \zeta^2+\zeta+1=0
\]
satisfy the equations and have $\beta\ne0$. If $\beta=0$, then $uy=1$, and the remaining minor is
$a_0b_2+a_2b_0=1+u$. Indeed the elliptic equation gives $1+u=x^3u^2$ under $uy=1$. Vanishing would force $u=1$ and $x=0$, contrary to $u=x^{q/2-2}$ and $q\ge6$.

At infinity the leading pole orders of $A,B$ have opposite parities, $2q-1$ and $2q$, and their leading coefficients are one. They span in the local Frobenius basis $(t^{-2q},t^{-2q+1})$. There is no rank drop there. At $Q_\lambda$, for $\lambda=\zeta,\zeta^2$, the function $x+y$ has a simple zero since $\partial(x+y)=1+x^2=\lambda\ne0$. Its factor $xu$ is a unit. Therefore the determinantal ideal of the three sections is the maximal ideal at each of these two points, with length one.

Let $K=\ker\Psi$. Away from the two points, trace-zero endomorphisms give all its sections by
\[
                \End^0(V)\longrightarrow K,\qquad
                T\longmapsto(Tf_0,Tf_1,Tf_2).
\]
This follows by using two of the $f_i$ as a basis: the single wedge condition on those two is precisely trace zero, and the other equations determine the third image. At a rank-drop point, elementary changes of bases put the triple in the form $(te_1,e_2,0)$ for a uniformizer $t$. In these coordinates
\[
 \Psi\bigl((a,b),(c,d),(e,f)\bigr)=(a+td,tf,e).
\]
Its kernel consists of $((td,b),(c,d),0)$. Trace-zero endomorphisms give the submodule in which $t$ divides $b$. Hence both $K/\End^0(V)$ and $\operatorname{coker}\Psi$ have length one at each of the two points. In particular $\deg K=2$, since $\deg\End^0(V)=0$.

The trace pairing identifies $(\End^0 V)^\vee$ with $\End(V)/\OO_E\operatorname{id}$. The local calculation says that its class lies in $K^\vee$ exactly when its representative preserves the image line of the three sections at the rank drop. In detail, the added local direction is $t^{-1}$ times the nilpotent sending the missing basis direction to the image line. The trace pairing extends integrally exactly when it annihilates that nilpotent modulo $t$, which is the stated preservation condition.

At $Q_\lambda$, row $a$ vanishes, so the image line is $kx$ in the rational basis $(1,x)$. By Lemma~\ref{lem:ss-end}, every global class is $\alpha[\operatorname{mult}_x]+\beta[\partial]$. The matrices \eqref{eq:ss-endmat} show that preservation imposes $\alpha\lambda+\beta=0$. The two conditions for $\lambda=\zeta,\zeta^2$ are independent. Thus $H^0(K^\vee)=0$, and duality and Riemann--Roch on the elliptic curve give $h^1(K)=0$ and $h^0(K)=\deg K=2$.

The bundles $V$ and $\det V$ have no $H^1$, by projection formula or positive-degree elliptic cohomology. Splitting the sequence for $\Psi$ at its image now identifies the cokernel of $H^0(\Psi)$ with the two-dimensional space of values of its torsion cokernel. It follows that $\rk H^0(\Psi)=6q-2$.

To check stable rank, evaluate kernel vectors in this cokernel. At $Q_\lambda$ one has $b=(0,\lambda,\lambda^2)$, so the cokernel functional on an output triple is
\[
        \ell_\lambda(U,V,W)=\lambda^2U(Q_\lambda)+\lambda V(Q_\lambda).
\]
The kernel contains $s_q,\partial_1s_q$ by Lemma~\ref{lem:upper}. At $Q_\lambda$, their triples of values are
\[
                (0,1,\lambda^2),\qquad(0,\lambda^2,\lambda),
\]
respectively. This is obtained by substituting $x=y=\lambda$ and differentiating the displayed $A,B,C$; the divisibility $6\mid q$ gives $x^{q/2}=1$. Their kernel--cokernel matrix is consequently
\[
 \begin{pmatrix}\ell_\zeta(s_q)&\ell_\zeta(\partial_1s_q)\\
 \ell_{\zeta^2}(s_q)&\ell_{\zeta^2}(\partial_1s_q)\end{pmatrix}
 =\begin{pmatrix}\zeta&1\\\zeta^2&1\end{pmatrix},\qquad\det=1.
\]
Thus the kernel has dimension two, meets the image trivially, and first rank already equals stable rank. Since the entire coefficient matrix is over $\FF_2$, the twist in \eqref{eq:twist} leaves it unchanged. The first two iterate ranks are both $6q-2$. The nonempty rank open meets the smooth open, and Lemma~\ref{lem:upper} gives the matching upper bound.
\end{proof}

\begin{lemma}[The supersingular bidegree $(6,3)$]\label{lem:ss-six}
A general smooth member of $|N(6,3)|$ on $E\times E$ has genus nineteen and full and relative stable rank sixteen.
\end{lemma}
\begin{proof}
Take the section
\[
 s=z+xw+x^2+x^2w+x^3+yz+xyz.
\]
Use the bases $(1,x,x^2,x^3,y,xy)$ and $(1,z,w)$ and tensor index $3i+j$. The one-variable Cartier rule is
\begin{align*}
 C(x^r dx)&=\begin{cases}0&r\text{ even},\\x^{(r-1)/2}dx&r\text{ odd},\end{cases}\\
 C(x^r y\,dx)&=\begin{cases}x^{r/2+1}dx&r\text{ even},\\x^{(r-1)/2}y\,dx&r\text{ odd}.
 \end{cases}
\end{align*}
It follows at once from $y=y^2+x^3$ and the coefficient rule. Multiplying by $s$, reducing $y^2=y+x^3$ and $w^2=w+z^3$, and applying this rule in both factors gives the eighteen columns of its Cartier matrix. Encoding a binary column $(v_0,\ldots,v_{17})$ by $\sum_i2^iv_i$, they are
\[
\begin{gathered}
4106,12,16432,4177,40,16660,32848,96,131456,\\
33416,320,133280,12424,49472,82080,65664,164608,66560.
\end{gathered}
\]
The finite rank certificate is recorded explicitly in Appendix~\ref{app:cert}: a specified sixteen-by-sixteen minor of the square is one. Lemma~\ref{lem:upper} bounds the first rank by sixteen, so both ranks equal sixteen and the image is stable. Openness, very ampleness and Lemma~\ref{lem:cohom} finish the proof.
\end{proof}

\section{Every even supersingular thin system}\label{sec:ss-all-even}

We now give a single section for every even first degree. Its Cartier kernel is a fixed rank-three bundle, and two residue functionals show that the image is already stable. This removes any divisibility-by-three condition from the supersingular thin theorem.

\begin{proposition}[An explicit section in every even degree]\label{prop:ss-all-even}
On the model $E/\FF_2$ given by $y^2+y=x^3$, write $z,w$ for the coordinates on a second copy. Let $m=2t\ge6$, and put
\[
 f=\begin{cases}
 x^{t/2},&t\text{ even},\\
 x^{(t-3)/2}y,&t\text{ odd},
 \end{cases}
 \qquad
 a=x+f^2,\qquad c=x^{t-2}y.
\]
The section
\[
                         s=a+z+cw\in H^0(E\times E,N(m,3))
\]
has first and stable Cartier rank $3m-2$. After extension to any algebraically closed field of characteristic two, a general smooth member of $|N(m,3)|$ has genus $3m+1$ and full and relative $2$-rank $3m-2$.
\end{proposition}
\begin{proof}
We carry out the bundle and residue calculations over $\FF_2$; all the displayed functions are defined over this field. Write $\omega=dx$ and $C(h)=C_E(h\omega)/\omega$. Since $x$ is a $2$-basis,
\[
 C(u^2+xv^2)=v,\qquad
 \partial(u^2+xv^2)=v^2.
\]
In the second factor,
\[
 C'(1)=C'(z^2)=C'(w^2)=0,\quad
 C'(z)=1,\quad C'(w)=z,\quad C'(zw)=w.
\]
These identities follow by substituting $w^2+w=z^3$ into the same coefficient rule. Hence, identifying a section of $N(m,3)$ with a triple in $H^0(E,\OO_E(mO))^3$, multiplication followed by Cartier is
\begin{equation}\label{eq:ss-even-cartier}
 T_s(t_0,t_1,t_2)=
 \bigl(C(t_0+at_1),\ C(ct_0+at_2),\ C(ct_1+t_2)\bigr).
\end{equation}

Introduce
\[
 (g,h)=
 \begin{cases}
 (x^{t/2},\,x^{(t-2)/2}y),&t\text{ even},\\
 (x^{(t-3)/2}y,\,x^{(t+1)/2}),&t\text{ odd},
 \end{cases}
 \qquad H=fg+h.
\]
The elliptic equation gives
\begin{equation}\label{eq:ss-even-identities}
 c=h^2+xg^2,\qquad \partial a=1,\qquad
 \partial c=g^2,\qquad ag^2+c=H^2.
\end{equation}
The functions $f,g,h$ have respective pole orders $t,t,t+1$ at $O$, with leading coefficient one for the parameter $\tau=x/y$. Thus $a,c,H$ have respective pole orders $m,m-1,m$, and $s$ is a section of the asserted bundle.

\paragraph{The exact kernel bundle.}
Put $L_1=\OO_E(mO)$, and let $F:E\to E$ be Frobenius on the $\FF_2$ model. Formula~\eqref{eq:ss-even-cartier} defines an $\OO_E$-linear map
\[
                   T:(F_*L_1)^3\longrightarrow L_1^3.
\]
The source module structure is through squares. We claim that
\begin{equation}\label{eq:ss-even-kernel}
 0\longrightarrow F_*\OO_E\oplus\OO_E
 \xrightarrow{\ i\ }(F_*L_1)^3
 \xrightarrow{\ T\ }L_1^3\longrightarrow0,
 \qquad i(q,r)=qs+r^2\partial s
\end{equation}
is exact, where $\partial s=(1,0,g^2)$.

Over the function field, put $q=t_1$ for a vector in the kernel and subtract $qs$. The first and third equations in~\eqref{eq:ss-even-cartier} leave a vector $(u,0,v)$ with $C(u)=C(v)=0$, so $u=r^2$ and $v=r_2^2$. The middle equation is $rg+r_2=0$; hence the vector is exactly $qs+r^2\partial s$. Conversely, substitution proves that every such vector is in the kernel.

Away from $O$, regularity of the vector gives regularity of $q=t_1$ and then of $r^2=t_0+aq$, hence of $r$. At $O$, the coefficient matrix for $q,r^2$ in the first and third components, after multiplying them by $\tau^m$, is
\[
 \begin{pmatrix}
 \tau^ma&\tau^m\\
 \tau^mc&\tau^mg^2
 \end{pmatrix}.
\]
Its determinant is $\tau^{2m}H^2$, a unit. Thus membership in $L_1^3$ at $O$ is equivalent to regularity of $q,r$. This proves the kernel assertion as a statement of bundles.

The quotient by that kernel has rank three and degree $3m$: finite Frobenius pushforward preserves Euler characteristic, and on an elliptic curve Euler characteristic equals degree. The target $L_1^3$ also has degree $3m$. The induced generically injective map between these bundles consequently has torsion cokernel of length zero. This proves the surjectivity in~\eqref{eq:ss-even-kernel}.
Since $H^1(E,L_1)=0$, global sections give
\begin{equation}\label{eq:ss-even-global-kernel}
 \ker T_s=\langle s,\partial s\rangle,\qquad
 \operatorname{coker}T_s
       \simeq H^1(E,F_*\OO_E\oplus\OO_E).
\end{equation}
Both spaces have dimension two, so $\rk T_s=3m-2$.

\paragraph{The first residue functional.}
For an output triple $u=(u_0,u_1,u_2)$, set
\[
 \ell(u)=\operatorname{Res}_O
                \frac{gu_0+u_2}{H}\,\omega.
\]
It kills the image of $T_s$. Indeed, by~\eqref{eq:ss-even-identities},
\[
 \frac{g}{H}T_0(t)+\frac1H T_2(t)
        =C(t_1)+C\!\left(\frac{g^2t_0+t_2}{H^2}\right).
\]
The first differential has no possible pole away from $O$, so its residue there is zero by the global residue theorem. The argument of the second Cartier operator is regular at $O$, because its numerator has pole order at most $2m$ and its denominator has pole order $2m$. Cartier preserves regular differentials.

Put
\[
                       R=\frac{c+xg+fh}{H}.
\]
The terms $c/H$ and $xg/H$ are regular at $O$. The term $fh/H$ has a simple pole there with coefficient one. Since $\omega=(1+O(\tau))\,d\tau$, this gives
\begin{equation}\label{eq:ss-even-residue}
                   \operatorname{Res}_O(R\omega)=1.
\end{equation}
The identity $(ag+c)/H=f+R$ and the global residue theorem for $f\omega$ imply $\ell(s)=1$. On the other hand, $(g+g^2)/H$ is regular at $O$, so $\ell(\partial s)=0$.

\paragraph{The second residue functional.}
Project the connecting map from~\eqref{eq:ss-even-global-kernel} onto the second factor $H^1(E,\OO_E)$ and pair with $\omega$ by residue duality. Denote the resulting cokernel functional by $\eta$. The residue of a principal part times $\omega$ is well-defined: regular terms and rational functions with their only pole at $O$ have zero residue. It is nonzero on $\tau^{-1}$, and hence identifies the one-dimensional group $H^1(E,\OO_E)$ with $\FF_2$.

We calculate $\eta(\partial s)$ by local preimages. On $E\setminus\{O\}$ take
\[
                  t_U=(x,0,h^2+ag^4).
\]
Formula~\eqref{eq:ss-even-cartier} gives $T_s(t_U)=\partial s$: use $C(x)=C(a)=1$, $C(h^2)=0$, and
\[
 C(cx)=h,\qquad
 C(ah^2+a^2g^4)=h.
\]
Near $O$ put
\[
 q=\frac{ag^4+c}{H^2},\qquad r=H+R,\qquad t_O=(0,q,0).
\]
The pole order of $q$ is at most $m$. After excluding the other zeros of $H$, the vector $t_O$ is therefore a local section of $L_1^3$. Direct substitution in the definitions gives
\begin{equation}\label{eq:ss-even-square}
                              r^2=x+aq.
\end{equation}
To verify this without an implicit cancellation, multiply by $H^2$. The two sides become equal because
\[
 H^4+c^2+x^2g^2+f^2h^2
   =xH^2+a^2g^4+ac
   =f^4g^4+x^2g^4+x^2g^2+f^2h^2.
\]
Furthermore,
\[
 qc+r^2g^2=q(c+ag^2)+xg^2
           =ag^4+c+xg^2=ag^4+h^2.
\]
Consequently $t_U+t_O=qs+r^2\partial s$, so $t_O$ is also a preimage of $\partial s$ and the connecting cocycle is represented by $(q,r)$. Its second residue is
\[
 \eta(\partial s)=\operatorname{Res}_O(r\omega)
 =\operatorname{Res}_O(H\omega)+\operatorname{Res}_O(R\omega)=1.
\]
Here $H$ has no pole away from $O$, and~\eqref{eq:ss-even-residue} supplies the second term.

The two cokernel functionals on the kernel basis $s,\partial s$ therefore have matrix
\[
                         \begin{pmatrix}1&0\\ *&1\end{pmatrix}.
\]
It is invertible. Kernel and image intersect trivially and have complementary dimensions, so the image is stable and the first rank already equals stable rank.

Finally, all coefficients of the section and its Cartier matrix are in $\FF_2$. After extension to $k$, absolute Cartier acts by that same matrix followed by inverse Frobenius on coefficients. The stable rank is therefore unchanged; no untwisted matrix power has replaced an arbitrary semilinear iterate. The resulting nonempty iterated-rank open meets the smooth open of the very ample complete system. Lemmas~\ref{lem:cohom} and~\ref{lem:upper} give connectedness, genus, and the matching bound $3m-2$. This proves the proposition.
\end{proof}


\section{Nodal images under one-factor isogenies}\label{sec:norm}

We now transfer the starting ranks to larger rectangular systems. Acting in only one factor makes the necessary point and tangent separation explicit, even when the \'etale degree is even.

\begin{lemma}[Elliptic \'etale isogenies with prescribed target]\label{lem:isogeny}
Let $E/k$ be an elliptic curve. For every odd $e>0$ there is an \'etale isogeny $E'\to E$ of degree $e$. If $E$ is ordinary, this holds for every $e>0$. The source has the same ordinary or supersingular type as the target.
\end{lemma}
\begin{proof}
For odd $e$, choose a cyclic subgroup of order $e$ in $E[e](k)$. The quotient isogeny has degree $e$, and its dual is an isogeny onto $E$ with the same degree. Since $e$ is prime to the characteristic, both are \'etale. For the standard degree and duality statements see \citet[Theorem~7.5 and Corollary~7.17, pp.~100, 104]{EGM}.

For an ordinary curve, Verschiebung $E^{(2)}\to E$ is \'etale of degree two. Indeed $[2]=VF$, Frobenius has degree two, and the ordinary curve has two geometric two-torsion points; therefore the separable degree of $V$ is its full degree. Iterating and composing with an odd-degree isogeny proves the claim for $e=2^ve_0$. These factorization and degree statements are \citet[Propositions~5.20--5.22, pp.~80--81]{EGM}. Finally $p$-rank is isogeny invariant, for example by the invariance of the rational \'etale Tate module, so the source has the asserted type.
\end{proof}

\begin{lemma}[One-factor nodal image]\label{lem:axis}
Let $m,n$ be positive multiples of three, let $E_1',E_2$ be elliptic curves, and let $\phi:E_1'\to E_1$ be a finite \'etale isogeny of degree $e$. For a general smooth curve $C'\in|N(m,n)|$ on $E_1'\times E_2$, the map $\phi\times1$ sends $C'$ birationally onto an integral nodal curve. Its numerical bidegree is $(m,en)$ and it has exactly $mn(e-1)$ nodes. Translation places the image in the exact system $|N(m,en)|$. The same statement with the factors exchanged sends bidegree $(m,n)$ to $(em,n)$.
\end{lemma}
\begin{proof}
Over $k$, the kernel is a finite constant group. Fix a nonzero kernel translation $g$, and a point $P=(p,q)$, so the pair is $(p,q),(p+g,q)$. The two value conditions are independent, because degree at least three on an elliptic curve separates two points. If $m\ge6$, the degree-$m$ bundle separates the first jets at $p,p+g$: after subtracting $2p+2(p+g)$ its degree is positive, so $H^1$ vanishes. Tensor with a section nonzero at $q$ to obtain arbitrary values and first-direction derivatives at both points. A section on the second factor vanishing at $q$ with nonzero first derivative, tensored with arbitrary first-factor values, gives the other two derivatives independently. Thus all six surface jet coordinates are available.

If $m=3$, the four first-factor jet coordinates have rank three. The kernel of their evaluation is zero, since a degree-three section cannot vanish to order at least two at each of two points. The two values are independent; after setting them to zero, the pair of first-direction derivatives spans a nonzero line, say $(a,b)$ with $(a,b)\ne(0,0)$. The two second-direction derivatives remain independent through tensor products as above. Writing the first-direction derivatives as $t(a,b)$, tangency after the two values are imposed is the nonzero equation $t(av+bu)=0$. In both cases the two vanishing values and tangency impose codimension three in the section space. Allowing $P$ to range over the two-dimensional surface gives a proper closed bad locus in the projective system.

For distinct nonzero kernel elements $g,h$, consider the three points with first coordinates $p,p+g,p+h$ and common second coordinate $q$. If $m\ge6$, their values are independent, since the residual degree after subtracting the three points is positive. If $m=3$, their value rank is three except when
\[
                       3p+g+h=0
\]
in the elliptic group, when it is two. The exceptional first coordinates form a finite set because $[3]$ is finite; the corresponding surface locus has dimension one. Thus the ordinary triple incidence has three conditions over dimension two, and the exceptional incidence has two conditions over dimension one. Both give proper bad loci. There are finitely many pairs and triples of kernel elements. A general smooth source curve consequently has no tangential pair and no triple identification.

A smooth ample source curve is connected, by Lemma~\ref{lem:cohom}, and therefore integral. If a nonzero kernel translation preserved it, all its pairs would be tangential, which has been excluded. Its image is thus birational, and because the ambient map is \'etale, the only remaining identifications are transverse pairs. The image is integral and nodal.

Pushforward of a vertical fibre has degree one onto its image, whereas pushforward of a horizontal fibre has degree $e$. This gives bidegree $(m,en)$. The image has arithmetic genus $emn+1$ and normalization genus $mn+1$, so the number of nodes is $mn(e-1)$. Its line bundle differs from $N(m,en)$ by an element of $\Pic^0(E_1\times E_2)$. The polarization map of an ample line bundle is an isogeny and hence is surjective on geometric points; see \citet[Proposition~6.13, pp.~93--94]{EGM}. A target translation therefore gives the exact line bundle. Translation preserves the singularities and the stable ranks. The proof with factors exchanged is identical.
\end{proof}

\begin{corollary}[Transfer of a sharp defect]\label{cor:transfer}
Suppose a nonempty smooth open in $|N(m,n)|$ has full curve defect $\delta$. Under either isogeny in Lemma~\ref{lem:axis}, a general smooth target curve has defect at most $\delta$. This holds with any permitted \'etale degree, including even degree on an ordinary factor.
\end{corollary}
\begin{proof}
Choose the source simultaneously in the given open and in the nodal-image open. By Lemma~\ref{lem:nodal}, the image has the same defect as its normalization: each node adds one both to the arithmetic genus and to the stable rank. Translate it to the exact target system. The open condition on a sufficiently high Frobenius iterate then meets the smooth target open, proving the bound.
\end{proof}

\section{Assembly of the rectangular formula}\label{sec:assembly}

We now combine the cubic, thin-system, and nodal calculations, and account for every parity case.

\begin{lemma}[Two ordinary factors]\label{lem:ordinary-grid}
If both elliptic factors are ordinary, a general smooth member of $|N(m,n)|$, for $m,n\ge3$, is ordinary.
\end{lemma}
\begin{proof}
Choose reduced divisors of degrees $m,n$ in the two systems $|\OO(mO)|,|\OO(nO)|$. The union of the corresponding $m$ vertical fibres and $n$ horizontal fibres is a connected nodal divisor in $|N(m,n)|$. All its normalization components are ordinary elliptic curves, so Lemma~\ref{lem:nodal} gives defect zero. The rank open meets the smooth open and gives ordinary smooth members. Since the stable rank cannot exceed the genus, this is the generic rank.
\end{proof}

\begin{lemma}[The odd side of a mixed rectangle]\label{lem:mixed-odd}
On an ordinary first factor and supersingular second factor, a general smooth member of $|N(3A,3B)|$ has defect one whenever $A$ is odd, with no condition on $B>0$.
\end{lemma}
\begin{proof}
Choose an \'etale isogeny $E_1'\to E_1$ of degree $B$ and an \'etale isogeny $E_2'\to E_2$ of degree $A$, using Lemma~\ref{lem:isogeny}. On $E_1'\times E_2'$, Proposition~\ref{prop:cubic} gives a smooth open in bidegree $(3,3)$ of defect one. First act on the first factor. Corollary~\ref{cor:transfer} gives a smooth curve of defect at most one in bidegree $(3,3B)$ on $E_1\times E_2'$. The parity bound is one, so a nonempty sharp open exists there. Now act on the second factor; this gives defect at most one in bidegree $(3A,3B)$. The ambient supersingular cohomology, or Lemma~\ref{lem:upper}, prevents smaller defect. Thus equality holds. The order of the two transfers is important: a degree $B$ isogeny on the first factor multiplies the second degree.
\end{proof}

\begin{lemma}[Odd supersingular rectangles]\label{lem:ss-odd}
On two supersingular factors, a general smooth curve in $|N(3A,3B)|$ has defect two when $A,B$ are odd.
\end{lemma}
\begin{proof}
Choose the two odd-degree \'etale isogenies onto the prescribed factors as in the preceding proof. All supersingular source curves are isomorphic to the model of Section~\ref{sec:cubic}. Begin with its defect-two cubic and transfer successively in the two factors. The defect is at most two, while Lemma~\ref{lem:upper}, together with the zero ambient stable rank, gives defect at least two. Hence the two are equal.
\end{proof}

\begin{proof}[Proof of Theorem~\ref{thm:thin}(2)]
Every supersingular elliptic curve over $k$ is isomorphic, with its origin, to the model $y^2+y=x^3$. Proposition~\ref{prop:ss-all-even} therefore applies to the prescribed two factors for every even $m\ge6$. Its sharp rank open meets the smooth open; Lemmas~\ref{lem:cohom} and~\ref{lem:upper} give genus $3m+1$ and full and relative rank $3m-2$.
\end{proof}

\begin{proof}[Proof of Theorem~\ref{thm:main}]
Put $A_0=ad$, $B_0=bd$, so the target bidegree is $(3A_0,3B_0)$. If both factors are ordinary, Lemma~\ref{lem:ordinary-grid} gives $f_C=9abd^2+1$. Subtracting ambient rank two gives \eqref{eq:main}.

Suppose the first factor is ordinary and the second is supersingular. If $A_0$ is odd, Lemma~\ref{lem:mixed-odd} gives defect one. If $A_0$ is even, take an ordinary \'etale source isogeny of degree $B_0$ onto the first factor. Theorem~\ref{thm:thin}(1) gives a sharp smooth source in bidegree $(3A_0,3)$, of defect two, on the product of that source curve and the given supersingular factor. Acting in the ordinary factor transfers it to bidegree $(3A_0,3B_0)$. Lemma~\ref{lem:upper} gives the matching defect two. Subtracting ambient rank one proves the mixed formula. Exchanging factors proves the other mixed ordering. This proof also includes the small degree-one systems with $B_0=1$.

Suppose finally that both factors are supersingular. If both $A_0,B_0$ are odd, Lemma~\ref{lem:ss-odd} applies. If exactly one is even, exchange factors if needed so that $A_0$ is even. Theorem~\ref{thm:thin}(2) gives defect three in bidegree $(3A_0,3)$. An odd-degree isogeny of degree $B_0$ on the first factor transfers it to $(3A_0,3B_0)$, and the parity bound gives equality.

If both $A_0,B_0$ are even, decompose the target class as
\[
               (3A_0,3B_0)=(3,3)+(3(A_0-1),3(B_0-1)).
\]
Both summands have odd coefficients, so Lemma~\ref{lem:ss-odd} gives smooth members of defect two in each. Their very ample systems allow them to be chosen transverse while remaining in their sharp opens: choose the first member in its open, and then use the basepoint-free restriction of the second system to avoid tangencies. Their positive intersection number makes the nodal union connected. Its defect is four by Lemma~\ref{lem:nodal}. The rank open intersects the smooth open in the target system, giving defect at most four, and Lemma~\ref{lem:upper} gives at least four. Thus all supersingular cases yield \eqref{eq:main}.

The dimensions and genera follow from Lemma~\ref{lem:cohom}. All arguments are geometric, so extension from a perfect ground field to its algebraic closure gives the stated geometric generic ranks. Finally, if $\mathcal I_A$ is the ideal sheaf in the embedding by $L$, the standard regularity vanishing $H^1(\mathcal I_A(d))=0$ holds for $d\ge\operatorname{reg}(\mathcal I_A)-1$. This is the definition and persistence of regularity in \citet[Definition~33.35.7 and Lemma~33.35.10]{Stacks}. The restriction map on sections is then surjective. Its inverse image of a nonempty rank open is a nonempty open in the ambient section parameter space, which proves the ambient assertion.
\end{proof}

\section{Scope beyond product polarizations}

The preceding theorem leaves no parity or valuation exception for the stated rectangular systems. It does not give the generic rank for an arbitrary inseparable polarization on a nonproduct abelian surface. For an ample line bundle $M$ on a mixed abelian surface $B$ of $2$-rank one, put $\lambda_M(x)=t_x^*M\otimes M^{-1}$ and define
\[
 a=\operatorname{length}_{\mathbf Z_2}\operatorname{coker}T_2(\lambda_M),
 \qquad b=v_2(\chi(M))-a.
\]
Here $T_2$ is the geometric Tate module and $v_2$ is the two-adic valuation. The unresolved primitive mixed case has $a>0$ and $b=0$. The missing sharpness assertion asks for the nonvanishing of the restricted Cartier determinant in that case. The ordinary elliptic evaluation isomorphisms used here need not exist on that surface, so the residue argument does not settle that condition. The theorem also concerns stable rank, not the complete Newton polygon, the group scheme of all two-torsion, or a monodromy image.

\appendix
\section{Finite Cartier certificates}\label{app:cert}

This appendix makes the finite matrix steps reproducible from their entries and states an additional uniform mixed bidegree-$(6,3)$ witness. The infinite families in the main text are proved by the residue and bundle arguments, not by these finite calculations.

For the supersingular matrix $C$ in Lemma~\ref{lem:ss-six}, form $C^2$ by binary matrix multiplication. Its minor on the zero-based row and column lists
\begin{align*}
 I&=(0,1,2,3,4,5,6,7,8,9,10,11,12,14,15,16),\\
 J&=(0,1,2,3,4,5,6,7,8,9,10,11,12,14,15,17)
\end{align*}
has determinant one. Here is a short elimination certificate. At columns $j=0,\ldots,15$, successively swap row $j$ with row
\[
                0,2,3,3,4,7,6,7,10,9,12,11,12,13,14,15,
\]
respectively, then add that row to every lower row having a one in column $j$. Every resulting pivot is one. This algorithm, together with the eighteen displayed columns, specifies each arithmetic operation of the rank certificate.

\begin{proposition}[A uniform mixed bidegree-$(6,3)$ witness]\label{prop:mixed-six}
On $E_c\times E_{\rm s}$ with $c\ne0$, the section
\[
       s=1+z+x(1+z+w)+xyz+x^3(z+w)
\]
has first and stable Cartier rank sixteen. In particular this witness works for every ordinary elliptic parameter.
\end{proposition}
\begin{proof}
Use the ordered first-factor basis $(1,x,y,x^2,xy,x^3)$ and second-factor basis $(1,z,w)$. Put $u=y+\sqrt c$. The first-factor pairing is
\begin{equation}\label{eq:beta6}
\begin{pmatrix}
1&0&x&x&u&0\\
0&x&u&0&x^2&x^2\\
x&u&y&x^2&0&xu\\
x&0&x^2&x^2&xu&0\\
u&x^2&0&xu&xy&x^3\\
0&x^2&xu&0&x^3&x^3
\end{pmatrix}.
\end{equation}
To verify all its entries without a table assumption, write the six basis elements as $A_i^2+xB_i^2$, with
\[
 A=(1,0,x,x,u,0),\qquad B=(0,1,u/x,0,x,x).
\]
The equation of $E_c$ verifies these decompositions, and the pairing entry is $A_iA_j+xB_iB_j$. Reduce it in the displayed basis; this is \eqref{eq:beta6}.

Let $T$ be the tensor contraction of \eqref{eq:beta6} with $\beta_{\rm s}$ and the eight coefficients of $s$. Explicitly, if $\beta_{ij}^r$ and $\gamma_{uv}^t$ denote the basis coefficients of the pairings, then
\[
 T_{3r+t,3j+v}=\sum_{i,u}s_{3i+u}\,\beta_{ij}^r\gamma_{uv}^t,
 \qquad s_\ell=1\ \text{for }\ell\in\{0,1,3,4,5,13,16,17\},
\]
and the other $s_\ell$ vanish. Form $M=(T^{[2]})^t$ and $S=MM^{[2]}$. The sixteen-by-sixteen minor of $S$ with lists
\begin{align*}
 I&=(0,7,2,3,8,5,6,9,11,1,10,14,13,15,4,12),\\
 J&=(0,1,2,3,4,5,6,7,9,10,11,12,13,14,15,16)
\end{align*}
has determinant $c^5$. For a verifiable elimination over $\FF_2(c)$, set
\[
 A_0=1+c+c^3,\quad B_0=1+c^2+c^4+c^5+c^6,
 \quad E_0=1+c+c^2,\quad D_0=E_0^2.
\]
Ordinary elimination in the listed order requires no row exchanges and has successive pivots
\[
 A_0,\ \frac c{A_0},\ cB_0,\ 1,\ \frac{c^4}{B_0},\ 1,
 \ \frac{D_0}{c^2},\ \frac1{cD_0},\ c,\ \frac{E_0}c,
 \ \frac{c^3}{E_0},\ \frac1c,\ c^2,\ \frac1c,\ \frac1c,\ 1.
\]
Their product is $c^5$. Since the original minor is a polynomial, this identity over the rational function field is a polynomial identity, valid even where an intermediate denominator vanishes. This also provides an independent expansion check: in characteristic two a determinant is its permanent, so it can be computed by the recurrence
\[
 \Delta_0(\varnothing)=1,\qquad
 \Delta_{r+1}(T)=\sum_{j\in T}\Delta_r(T\setminus\{j\})\,B_{r,j}
 \quad(|T|=r+1),
\]
where $B$ is the chosen minor and its rows are numbered from zero. The final value is the same monomial $c^5$.

The parity bound gives first rank at most sixteen. The nonzero second-iterate minor gives rank at least sixteen. Equality of the first two ranks proves the stable-rank assertion. This finite witness is consistent with, and is a special case of, the uniform residue theorem.
\end{proof}

\section*{Sources and method}
The proof uses Cartier theory, elliptic cohomology, residue duality and elementary finite-isogeny theory at the primary locations cited above. The nodal and stable-part arguments also occur in the companion treatment of abelian-surface polarizations; the arguments required here are written in full. The accompanying \emph{Rectangular Cartier certificate} reconstructs the monomial operators and proves the parameter minors by exact subset expansion over $\FF_2[c]$, independently of the earlier fraction-free and symbolic implementations. The \emph{Thin Cartier certificate} checks the new sections over the four-element field, applying inverse Frobenius to coefficients at every iterate. These finite checks corroborate the uniform bundle and residue proofs; they do not infer an infinite family from sampled degrees. The source and certificates are available at \url{\workbenchurl}; permanent records are at \url{\recordurl}.

The mathematics and text were produced by GPT-6 Astra Pro in ChatGPT Pro mode. The author made a statement-led partial reconstruction with earlier exposure to companion rank proofs, followed by full targeted checks of the printed arguments and primary inputs. Original certificates were read and rerun, and the finite computations were reconstructed by a separate implementation. The complete manuscript also received a separate expert reading by GPT-6 Astra Pro, which found no mathematical defect. Complete separate proof-unseen reconstructions of both even thin-system assertions were frozen and compared in full; their exact kernel and residue computations also yield the supersingular statement for every even degree at least six. The rectangular main theorem also has a complete separate reconstruction, using precisely the ordinary-surface theorem and the degree-one principal-polarization theorem as declared companion inputs. Its nodal-image and mixed ordinary-fibre constructions were compared in full. These are same-model checks, not human or formal verification. Reasoning-effort metadata is not exposed. No human has checked the mathematics. The new text and code are dedicated to the public domain under CC0 1.0.

\begin{thebibliography}{99}
\bibitem[Chang et~al.(2026)]{ChangEtAl}
Chang, K., Dragutinovi\'c, D., Groen, S., Lin, Y., Pacheco-Tallaj, N., \& Singhal, D. (2026). The $p$-rank stratification of the moduli space of double covers of a fixed elliptic curve. \emph{Nagoya Mathematical Journal, 261}, Article e26, 1--31. \url{https://doi.org/10.1017/nmj.2026.10107}.
\bibitem[Edixhoven et~al.(n.d.)]{EGM}
Edixhoven, B., van der Geer, G., \& Moonen, B. (n.d.). \emph{Abelian varieties: Preliminary version of the first chapters}. Author manuscript. \url{https://van-der-geer.nl/~gerard/AV.pdf}.
\bibitem[Grothendieck(n.d.)]{Grothendieck}
Grothendieck, A. (n.d.). \emph{Notes on crystalline cohomology}, cote~4, \S8.5.8, archive sheets~7--8. Institut Montpelli\'erain Alexander Grothendieck. \url{https://grothendieck.umontpellier.fr/4.pdf}.
\bibitem[Illusie(1979)]{Illusie}
Illusie, L. (1979). Complexe de de Rham--Witt et cohomologie cristalline. \emph{Annales scientifiques de l'\'Ecole normale sup\'erieure, 12}(4), 501--661. \url{https://doi.org/10.24033/asens.1374}.
\bibitem[Oda(1969)]{Oda}
Oda, T. (1969). The first de Rham cohomology group and Dieudonn\'e modules. \emph{Annales scientifiques de l'\'Ecole normale sup\'erieure, 2}(1), 63--135. \href{https://doi.org/10.24033/asens.1175}{DOI: 10.24033/asens.1175}.
\bibitem[Oda(1971)]{Oda1971}
Oda, T. (1971). Vector bundles on an elliptic curve. \emph{Nagoya Mathematical Journal, 43}, 41--72. \url{https://doi.org/10.1017/S0027763000014367}.
\bibitem[Stacks Project Authors(2026)]{Stacks}
Stacks Project Authors. (2026). \emph{The Stacks Project}, \S33.35, Definition~33.35.7 and Lemma~33.35.10. \url{https://stacks.math.columbia.edu/tag/08A2}.
\bibitem[Tate(1968)]{Tate}
Tate, J. (1968). Residues of differentials on curves. \emph{Annales scientifiques de l'\'Ecole normale sup\'erieure, 1}(1), 149--159. \url{https://doi.org/10.24033/asens.1162}.
\bibitem[Zou(2026)]{Zou}
Zou, H. (2026, September 29). \emph{Quasi-$F$-split primitive symplectic varieties in positive characteristic}. Author preprint. \url{https://www.math.uni-bielefeld.de/~hzou/papers/qfpsv/}.
\end{thebibliography}
\end{document}
